\section{Proof Appendix Scaffold}

\paragraph{Theorem (time-uniform gap certificate, caveated).}
For the pre-registered bounded sequence $d_i=(a_i-C_i+1)/2$, an anytime-valid
confidence sequence controls coverage under optional stopping. A certified
claim may be drafted only when the lower bound on $\Delta$ exceeds the declared
threshold and all evidence-eligibility gates pass.

\paragraph{Caveat label.}
This theorem is conditional on the declared item population, parser, provider
set, and inclusion rules. Bootstrap summaries are descriptive only.

\section{V6 Confidence-Sequence Bridge}

\paragraph{Bounded variable.}
For each certificate-eligible item, define the bounded variable
$d_i=(a_i-C_i+1)/2 \in [0,1]$, where $a_i$ is original-image correctness and
$C_i$ is visual decision-update consistency.

\paragraph{Estimand.}
The estimand is the intervention-consistency gap
$\Delta = \mathbb{E}[a_i] - \mathbb{E}[C_i] = 2\mathbb{E}[d_i]-1$ for the
declared, validity-gated item population.

\paragraph{Optional stopping.}
The native implementation in \texttt{certvic/metrics/anytime_cs.py} is intended
to supply an anytime-valid confidence sequence under the predeclared item order
and stopping rule. The bridge between implementation assumptions and theorem
assumptions is proof-required until formal references are added.

\paragraph{Implementation bridge.}
The reporting path must cite \texttt{certvic/metrics/anytime_cs.py},
\texttt{certvic/metrics/certification.py}, and
\texttt{certvic/validity/filter_scores.py}; it must not cite fake papers or
interpret bootstrap summaries as certification.
